\section{Kritische Reflexion des Projekts}
Die Entwicklung der WordClock verdeutlicht die Komplexität interdisziplinärer Systeme.  
Die Verbindung von Elektronik, Software, Netzwerkkommunikation und mechanischer Konstruktion erforderte eine strukturierte Vorgehensweise sowie eine sorgfältige Abstimmung der einzelnen Teilbereiche.  

Im Projektverlauf zeigte sich, dass zusätzliche Funktionen den Ressourcenbedarf des Systems erhöhen und bereits in der frühen Auslegung berücksichtigt werden müssen.  
Dies betrifft insbesondere Speicherbedarf, Spannungsversorgung, Einschaltverhalten und den Integrationsaufwand der verschiedenen Komponenten.  

Auch im mechanischen Bereich waren mehrere Anpassungen notwendig, um Stabilität, Montagefähigkeit und eine geeignete Lichtverteilung zu erreichen.  

Insgesamt konnte das Projekt erfolgreich umgesetzt werden.  
Die gewonnenen Erkenntnisse bilden eine Grundlage für zukünftige Entwicklungen vergleichbarer elektrotechnischer Systeme.